% Глава: 04-filesystem
% Статус: наполнено (первый прогон).
\chapter{Файловая система}
\label{ch:04-filesystem}

\section{Единое дерево каталогов}

В Unix/Linux все ресурсы представляются как файлы, а файловая
система --- как \textbf{единое дерево} с одним корневым каталогом
\fpath{/}. В отличие от Windows (диски \fpath{C:}, \fpath{D:}),
здесь несколько физических дисков и разделов \emph{подключаются}
(\emph{mount}) в отдельные точки этого же дерева.

Ваш домашний каталог --- это \fpath{/home/<login>} на сервере
\fpath{pglab.nsu.ru}; всё, что вы создаёте, живёт на диске вашей ВМ.

Базовое устройство описывает стандарт \emph{FHS}
(Filesystem Hierarchy Standard). Ключевые каталоги:

\begin{center}
\begin{tabular}{ll}
\toprule
\textbf{Каталог} & \textbf{Назначение} \\
\midrule
\fpath{/etc}        & Конфигурация системных программ \\
\fpath{/home}       & Домашние каталоги пользователей \\
\fpath{/root}       & Домашний каталог \cmd{root} \\
\fpath{/usr}        & Основные программы и библиотеки системы \\
\fpath{/usr/bin}    & Исполняемые файлы (часть \fpath{PATH}) \\
\fpath{/var}        & Изменяемые данные: логи, кэши, почта \\
\fpath{/var/log}    & Файлы журналов системы и программ \\
\fpath{/tmp}        & Временные файлы (очищаются при перезагрузке) \\
\fpath{/dev}        & Файлы-устройства (диски, терминалы и т.п.) \\
\fpath{/proc}       & Виртуальная ФС с информацией о процессах и ядре \\
\fpath{/bin},\ \fpath{/sbin} & Критически важные команды (часто ссылки на \fpath{/usr}) \\
\bottomrule
\end{tabular}
\end{center}

На рис.\ \ref{fig:fs-tree} --- упрощённое представление дерева.

\begin{figure}[htb]
  \centering
  \begin{tikzpicture}[level distance=1.4cm, sibling distance=2.4cm]
    \node[process] {/}
      child {node[data] {/etc}}
      child {node[data] {/home}}
      child {node[data] {/usr}}
      child {node[data] {/var}}
      child {node[data] {/tmp}}
      child {node[data] {/dev}}
      child {node[data] {/proc}};
  \end{tikzpicture}
  \caption{Упрощённое дерево файловой системы Linux.}
  \label{fig:fs-tree}
\end{figure}

\section{Пути и навигация}

\begin{itemize}
  \item \fpath{/} --- корневой каталог;
  \item \fpath{~} --- домашний каталог текущего
    пользователя;
  \item \fpath{.} --- текущий каталог;
  \item \fpath{..} --- родительский каталог.
\end{itemize}

\textbf{Абсолютный путь} начинается с \fpath{/}
(\fpath{/home/alice/docs}); \textbf{относительный} --- от текущего
каталога (\fpath{docs/notes.txt}).

Основа навигации:
\begin{minted}[bgcolor=gray!6]{bash}
pwd                 # показать текущий каталог по абсолютному пути
ls                  # список содержимого
cd ~/work           # перейти в каталог
cd /                # перейти в корень
cd ..               # на уровень вверх
cd -                # вернуться в предыдущий каталог
ls -la ~/           # подробный список, включая скрытые файлы
\end{minted}

\cmd{ls} показывает \emph{видимое}; флаг \flag{l} ---
подробно (права, владелец, размер, время), \flag{a} --- включая
скрытые (имена, начинающиеся с \cmd{.}), \flag{h} --- размеры
в человекочитаемом виде, \flag{R} --- рекурсивно.
Все флаги всегда можно найти в \cmd{man ls}.

\section{Операции с файлами и каталогами}

\begin{center}
\begin{tabular}{lp{0.75\linewidth}}
\toprule
\textbf{Команда} & \textbf{Назначение} \\
\midrule
\cmd{mkdir}      & создать каталог (\flag{p} --- вместе с родителями) \\
\cmd{touch}      & создать пустой файл / обновить метки времени \\
\cmd{cp}         & копировать (\flag{r} --- рекурсивно, \flag{i} --- спросить, \flag{v} --- показывать) \\
\cmd{mv}         & переименовать или переместить \\
\cmd{rm}         & удалить файл (\flag{r} --- каталог с содержимым, \flag{i} --- спросить) \\
\cmd{rmdir}      & удалить \emph{пустой} каталог \\
\bottomrule
\end{tabular}
\end{center}

Типовые сценарии:
\begin{minted}[bgcolor=gray!6]{bash}
mkdir -p work/a/b/c     # создать цепочку каталогов за один шаг
touch work/a/b/c/one.txt work/a/b/c/two.txt
cp -rv work backup/     # показать, что копируется
mv one.txt renamed.txt  # переименовать
rm -i file.txt          # удалить с подтверждением
rm -rf work/a           # удалить всё дерево, осторожно!
mkdir empty && rmdir empty
\end{minted}

\begin{warnbox}
\cmd{rm -rf} на абсолютном пути не имеет ``корзины'' и не спрашивает.
Опечатка в пути (особенно в \fpath{/}) может удалить много данных.
Перед \cmd{rm -rf} перечитайте команду.
\end{warnbox}

\section{Метаданные, типы и размеры}

Файл в Linux --- это \emph{inode} (модель данных файловой системы,
см.\ далее) плюс имя. Команда \cmd{stat} показывает метаданные:
владельца, группу, права, размер, время изменения.
\cmd{file} определяет \textbf{тип} содержимого независимо от
расширения (по сигнатуре). \cmd{du} и \cmd{df} --- место:
\begin{minted}[bgcolor=gray!6]{bash}
stat config/service.conf       # метаданные
file image.bin                 # что это за файл на самом деле
du -sh ~/work                  # сколько места занимает каталог
df -h                          # свободное место на файловой системе
\end{minted}

\section{Поиск файлов}

Основной инструмент поиска --- \cmd{find}:
\begin{minted}[bgcolor=gray!6]{bash}
find . -name '*.conf'                 # все файлы с расширением .conf
find /home -type f -name 'user*'      # файлы, чьё имя начинается с user
find . -iname '*.log'                 # поиск без учёта регистра
find . -type d -name 'work'           # только каталоги
\end{minted}

\cmd{find} ищет по имени, типу, времени изменения и многим другим
атрибутам (\cmd{man find} --- одна из самых длинных справочных страниц).
Отдельная утилита \cmd{locate} ищет мгновенно по готовому индексу
имён, но индекс обновляется периодически.

\section{Как это устроено под капотом}

Файловая система хранит данные блоками на устройстве, а
\textbf{inode} --- это запись с атрибутами файла: права, владелец,
размер, указатели на блоки данных. Имя файла живёт в каталоге
(каталог --- это тоже файл, список «имя $\to$ inode»).
Это объясняет две вещи:
\begin{itemize}
  \item жёсткая связь (\emph{hard link}) --- ещё одно имя
    на тот же inode; удаление имени не удаляет данные, пока
    есть другие имена;
  \item \cmd{mv} внутри одной ФС очень дёшев --- перемещается
    только имя (запись в каталоге), а данные на месте.
\end{itemize}

Механизм inode объясняет и то, почему данные на вашей ВМ
(см.\ \nameref{ch:02-linux-access}) сохраняются после
разрыва SSH: они записаны в блочное хранилище ВМ,
а сессия --- лишь процесс, работающий с этим хранилищем.

\section{Типичные ошибки}
\begin{itemize}
  \item \emph{No such file or directory} --- опечатка в пути или
    файл не в текущем каталоге (забудьте \fpath{./}).
  \item Скрытые файлы ``не видно'' в \cmd{ls} --- нужен \flag{a}.
  \item \cmd{rmdir} не удаляет каталог --- он не пуст; используйте
    \cmd{rm -r} (после проверки содержимого).
  \item \cmd{mv a b}, где \fpath{b} --- каталог, перемещает
    \fpath{a} \emph{внутрь} \fpath{b}, а не переименовывает.
  \item Ищите файл, а \cmd{locate} не находит --- индекс устарел;
    используйте \cmd{find}.
\end{itemize}

\section{Связь с остальным курсом}
Права и владельцы, описанные в метаданных, --- тема
\nameref{ch:05-permissions}. Поиск и текстовые утилиты для анализа
файлов --- \nameref{ch:07-text-pipelines}. Монтирование точек
(откуда берётся дерево из нескольких дисков) и контейнеры,
изолирующие файловую систему, --- \nameref{ch:19-docker-basics}
и \nameref{ch:23-docker-resources}.

\section{Задачи для самопроверки}
\begin{enumerate}
  \item Объясните разницу между абсолютным и относительным путём;
    чем отличается \fpath{~} от \fpath{.}.
  \item Одной командой создайте каталог с тремя вложенными уровнями
    и пустой файл в самом глубоком.
  \item Определите тип произвольного файла, не смотря на расширение.
    Как вы это сделали?
  \item Найдите все файлы с расширением \fpath{.txt} в домашнем
    каталоге и подсчитайте их количество.
  \item Объясните на примерах, почему \cmd{mv} внутри одной
    файловой системы не копирует данные.
\end{enumerate}

